% GENERATED FILE. Edit canonical lessons, then run npm run generate:books.
% canonical-source-hash: fnv1a64:8cbc023c73c95858
% canonical-lessons: ES-C44-repaso-articulos, ES-C44-repaso-ir-ida

\chapter{Review --- All Four Words for The}
\label{ch:review-articles}
\hlchaptermodality{81}

\begin{chapteropening}
\textbf{By the end of this chapter:} \emph{I can give el, la, los and las and state the plural rule including its consonant repair.}

Complete ES-C44-repaso-articulos: give the four articles in order and state the plural rule in one sentence.
\end{chapteropening}

\section[Practice]{Review — all four words for \textquotedblleft{}the\textquotedblright{}}

\label{lesson:ES-C44-repaso-articulos}



\begin{quote}
Nothing new here. Before the table: how many words does Spanish have for \textquotedblleft{}the\textquotedblright{}?
\end{quote}

\begin{grammarlens}[title={What you've built}]
Four. And you have now met every one of them separately, which is why they are allowed to appear together.

\noindent
\begin{tabularx}{\linewidth}{@{}>{\raggedright\arraybackslash}X>{\raggedright\arraybackslash}X>{\raggedright\arraybackslash}X@{}}
\toprule
\textbf{} & \textbf{singular} & \textbf{plural} \\
\midrule
\textbf{masculine} & \textbf{el} libro & \textbf{los} libros \\
\textbf{feminine} & \textbf{la} casa & \textbf{las} casas \\
\bottomrule
\end{tabularx}

English has one word here and Spanish has four, so it is easy to read this table as three extra things to carry. It is not. It is \textbf{two} questions you were already answering — \emph{which gender?} and \emph{how many?} — laid out in a grid.

The first chapter of this book taught the top row. The last two chapters taught the bottom one. Nothing in this table arrived without being built.
\end{grammarlens}

\begin{grammarlens}[title={Grammar Lens: the two plural rules}]
And underneath the grid, one rule with a repair:

\begin{quote}
\textbf{Add \emph{-s}. If the word ends in a consonant and the \emph{-s} would be unsayable, add \emph{-es} instead.}
\end{quote}

\begin{itemize}
\raggedright
  \item \emph{casa} $\to$ \emph{casas}, \emph{libro} $\to$ \emph{libros}, \emph{comida} $\to$ \emph{comidas} — vowels, plain \emph{-s}.
  \item \emph{usted} $\to$ \emph{ustedes}, \emph{español} $\to$ \emph{españoles} — consonants, \emph{-es}.
\end{itemize}

You held the first half from chapter one and the second half from \emph{ustedes} without knowing it had a name.
\end{grammarlens}

\subsection*{Your turn}
Say these aloud:

\begin{itemize}
\raggedright
  \item \textquotedblleft{}el, la, los, las\textquotedblright{} until the order is automatic
  \item \textquotedblleft{}el libro, los libros, la casa, las casas\textquotedblright{}
\end{itemize}

\emph{Twice through:} Until the grid is a reflex rather than a lookup.

\begin{tcolorbox}[breakable,colback=teal!4,colframe=teal!35!black,title={Before you move on}]
Four articles — say them. (\emph{\textbf{El, la, los, las}}.) Two questions decide which? (\textbf{Gender} and \textbf{number}.) Next: what happens to \emph{a} when there is more than one — and the answer is not what you would expect.
\end{tcolorbox}
\section[(review)]{Review — the half of \emph{ir} you never say out loud}

\label{lesson:ES-C44-repaso-ir-ida}



\begin{quote}
\textquotedblleft{}We go.\textquotedblright{} (\emph{Vamos}.) \textquotedblleft{}They go.\textquotedblright{} (\emph{Van}.)
\end{quote}

\begin{grammarlens}[title={What you've built}]
Three plural forms, and a \emph{v} in every one of them.

\textbf{We.} \emph{Vamos.}  ·  \textbf{You all.} \emph{Vais.}  ·  \textbf{They.} \emph{Van.}

Now say the infinitive after them: \emph{ir}. No \emph{v} anywhere in it.

That gap is the whole point. The forms you speak descend from \textbf{vādere}, \textquotedblleft{}to make one's way.\textquotedblright{} The name on the dictionary entry descends from \textbf{īre}. Two Latin verbs, one Spanish label — so the present tense is not irregular at all. It is perfectly regular, for a verb whose own infinitive Spanish stopped saying.
\end{grammarlens}

\begin{culture}
So what is left of \emph{īre}? Almost nothing. The infinitive \emph{ir}, and one door you already know how to open.

Remember what \emph{comer} did when you took its ending off and put \textbf{-ida} on instead:

\begin{quote}
\emph{com}er $\to$ la \emph{com}ida — the eaten thing.
\end{quote}

Do the same to \emph{ir} and you get \textbf{la ida} — the going, the outward journey. Nobody had to teach you that word. The ending built it, and it built it on the one syllable \emph{vādere} never took over.

That is worth more than either word. An ending that turns a verb into the result of the verb keeps working on stems the spoken language has abandoned.
\end{culture}

\subsection*{Your turn}
Say these aloud:

\begin{itemize}
\raggedright
  \item \emph{vamos}, \emph{vais}, \emph{van} — then \emph{ir}, and hear the letter go missing
  \item \emph{la comida}, \emph{la ida} — one ending, two verbs
\end{itemize}

\emph{Twice through:}

\begin{tcolorbox}[breakable,colback=teal!4,colframe=teal!35!black,title={Before you move on}]
\textquotedblleft{}You all go\textquotedblright{}? (\emph{\textbf{Vais}}.) Which Latin verb do \emph{vamos}, \emph{vais} and \emph{van} belong to? (\emph{\textbf{Vādere}}.) And which one does \emph{la ida} reach? (\emph{\textbf{Īre}} — the half of \emph{ir} that stopped being spoken.)
\end{tcolorbox}
